\section{Sufficient geometry, counterexamples, and open problems}\label{art:verification}

The convergence conditions of Sections~\ref{art:global}--\ref{art:local} are
sufficient rather than necessary: Chapter~5 of the companion gives a
one-dimensional convex example with accumulating kinks on which every admissible
full-BFGS run converges without admitting a uniform sublevel-band certificate.
Complementing the positive theorems, the counterexamples of
Section~\ref{art:failure} explain why neither convexity, coercivity, nor smooth
or strongly convex branch formulas can replace those geometric conditions: when
$0\in\aff\partial F\setminus\conv\partial F$ eliminates equal-rate descent
directions, unmodified full BFGS from $H_0=I$ can enter a self-similar cycle and
converge to a nonstationary point.

\subsection{Why the counterexamples escape the positive theory}

Each counterexample violates a specific hypothesis of the positive theorems:
\begin{enumerate}
\item \emph{Global homogeneity vs.\ affine offsets.} The three-slope cone
$G(z)=\max_{0\le i\le2}g_i^Tz$ is positively homogeneous of degree one, but
satisfies $G(t(-1,16)^T)=-t$ for $t\ge0$ and so fails coercivity; its four-piece
coercive cap $\overline G$ in Theorem~\ref{art:failure_polyhedral} includes the
affine offset $v-17$, breaking homogeneity around the minimizer $(-1,16)$ and
leaving $(0,0)$ as a nonstationary corner with $0\in\aff\partial\overline G(0)\setminus\conv\partial\overline G(0)$.
\item \emph{Finite local polyhedral partitions vs.\ accumulating smooth kinks.}
Near $(-1,16)$, $\overline G$ is locally a three-piece polyhedral cone centered
at $(-1,16)$ and therefore converges locally by Theorem~\ref{art:polyhedral}.
By contrast, in Theorems~\ref{art:local-failure} and~\ref{art:local-failure-lower},
nonlinear scalar splines cause nonstationary three-slope corners $\bar x_\ell\to0$
to accumulate at the minimizer $0$: no neighborhood of $0$ admits a finite
polyhedral partition with smooth extensions (since Theorem~\ref{art:polyhedral}
would otherwise force convergence once $F(x_0^{(\ell)})\le\varepsilon$),
demonstrating that finitely many smooth formulas in $\max_if_i$ do not imply a
finite local polyhedral cell structure.
\item \emph{Active-gradient multiplier nondegeneracy.} Both local smooth-branch
counterexamples are planar four-branch maxima, so we can test the multiplier
hypotheses of Corollary~\ref{art:planar-max} directly at $x_*=0$. In
\eqref{art:upper_objective}, all four branches are active at $0$ with gradients
$0,0,0,e_2$, where $e_2=(0,1)^T$; in~\eqref{art:lower_objective}, the active
gradients at $0$ are again $0,0,0,e_2$ before whitening and $0,0,0,A^Te_2\ne0$
after $x=Az$ with $A=(H_0^{\rm seed})^{1/2}$. In both cases, the repeated zero
gradients violate affine independence, and any convex combination placing
positive weight $\lambda_3>0$ on the fourth branch equals $\lambda_3e_2\ne0$ (or
$\lambda_3A^Te_2\ne0$), violating strict complementarity $\lambda_i^*>0$ as
well; in~\eqref{art:lower_objective}, three branches also lose $C^2$ regularity
at $0$.
\end{enumerate}

A one-line secant identity also clarifies why uniform strong convexity cannot by
itself prevent BFGS from converging to a nonstationary point. On $Q_\mu(x)=\tfrac\mu2\norm{x}^2+G(x)$,
if consecutive differentiable iterates $x$ and $x+s$ activate slopes $g_i$ and
$g_j$, then the secant vector is
\begin{equation}\label{art:jump-secant}
 y=\bigl(\mu(x+s)+g_j\bigr)-\bigl(\mu x+g_i\bigr)
  =\mu s+(g_j-g_i).
\end{equation}
Along the perturbed trajectory of Proposition~\ref{art:quadratic_failure},
$\norm{s_k}\to0$ while $g_0,g_1,g_2$ continue to cycle: the strongly convex
contribution $\mu s_k$ vanishes with the step length, leaving the $O(1)$ jump
$g_j-g_i$ to dominate the rank-two BFGS update.

\subsection{Open problems}

Our results point directly to four open problems in the theory of quasi-Newton
methods for nonsmooth optimization:
\begin{enumerate}
\item \emph{Uniformly strongly convex $C^2$ branch maxima.} Globally,
$Q_\mu(x)=\tfrac\mu2\norm{x}^2+G(x)$ is already a maximum of three $\mu$-strongly
convex $C^\infty$ quadratic functions on which full BFGS can converge to the
nonstationary origin (Proposition~\ref{art:quadratic_failure}). Locally near the
minimizer, however, Table~\ref{art:regularity-map} leaves a sharp regularity
gap: Theorem~\ref{art:local-failure} achieves $C^{2,1}$ branches without a
positive lower Hessian bound, Theorem~\ref{art:local-failure-lower} achieves
uniformly strongly convex $C^{1,1}$ branches that lose $C^2$ regularity at $0$,
and Corollary~\ref{art:planar-max} proves local convergence in $\R^2$ when the
active gradients at $x_*$ are nondegenerate. Does unmodified full BFGS converge
locally from every differentiable start near the minimizer of a maximum of
finitely many uniformly strongly convex $C^2$ functions when active-gradient
nondegeneracy is dropped?
\item \emph{Moving comparison fields in dimension $n\ge3$.} The radial theorems
(Theorems~\ref{art:centered}, \ref{art:homogeneous}, and~\ref{art:polyhedral})
hold in arbitrary dimension $\R^n$, whereas our curved-interface analysis
(Theorem~\ref{art:planar-cluster} and Corollary~\ref{art:planar-max}) uses the
one-dimensionality of level-set tangents in $\R^2$ to cancel normal components
of $\partial_\theta(hn+h_\theta T)$ and charge field rotation to functional
decrease. Extending moving comparison fields, or the higher-rank comparison tensors
$W_\eta(x)$ already permitted by Theorem~\ref{art:certificate}, to nondegenerate
smooth maxima with curved switching manifolds in $\R^n$ ($n\ge3$) is a natural
next step.
\item \emph{Convergence rates via principal minors and comparison energies.}
Theorem~\ref{art:reciprocal-growth} and our comparison-energy inequalities give
asymptotic convergence without conditioning assumptions (finite-horizon
estimates are recorded in the companion). Determining sharp non-asymptotic or
R-linear convergence rates for structured nonsmooth families such as
$F(u,v)=u^2+|v|$ or coercive polyhedral norms remains largely open.
\item \emph{Positive-measure vs.\ zero-measure failure on convex functions
(Lewis--Overton Challenge~7.1).} On \emph{nonconvex} coercive semialgebraic
objectives, Lai~\preprintcite*{NEG_LAI26} recently constructed a nonempty open
set of initial states $(x_0,H_0)$ converging to a non-Clarke-stationary point by
using negative transverse curvature to contract perturbations onto a homogeneous
cycle. Under convexity, however, Theorem~\ref{art:homogeneous} rules out global
failure on every coercive positively homogeneous objective, and the failing
starts exhibited in Theorems~\ref{art:failure_polyhedral}, \ref{art:local-failure},
and~\ref{art:local-failure-lower} consist of a single point $z_0$ or discrete
sequences $x_0^{(\ell)}$ (or lie on a two-dimensional stable manifold in the
normalized state space of Appendix~\ref{art:lower_curvature}). Thus our
deterministic counterexamples do not contradict empirical success from typical
initializations: under absolutely continuous random initialization as in
Challenge~7.1~\cite{LO}, can a \emph{convex} nonsmooth function admit a
positive-measure set of initial states $(x_0,H_0)$ on which full BFGS converges
to nonstationary points, or are all convex nonstationary basins of measure zero?
\end{enumerate}

\subsection{Formal verification and reproducibility}

All results and rational interval certificates in this article, together with
the extended noncompact and quantitative theory in the accompanying companion
monograph, have been formalized and verified in the Lean~4 theorem prover on top
of \textsf{Mathlib}. Every verified declaration has been audited to depend only
on the three standard foundational axioms of Lean's classical logic
(\texttt{propext}, \texttt{Classical.choice}, and \texttt{Quot.sound}), with the
core library and companion modules additionally re-verified through an
independent \texttt{lean4checker} kernel replay. The pinned formal sources,
companion monograph, and reproducible audit scripts are publicly available at
\url{https://github.com/yxie1992/bfgs-nonsmooth}.

\paragraph{Acknowledgements.}
Large language models assisted theorem development, proof reconstruction, Lean~4
formalization, literature searches, software tooling, and manuscript preparation;
the companion report documents this workflow and the formal verification audits
in detail.

